% META: {"id": "TEXP-ARI-EX-013", "chapitre": "TEXP-ARITHMETIQUE", "type_objet": "exercice", "capacites": ["TEXP-ARI-C2"], "capacites_codes": ["C2"], "methodes": ["M2"], "parcours": 3, "competences": ["raisonner"], "duree_min": 18, "mode_creation": "ex_nihilo", "sources_inspiration": [], "fichier_tex": "chapitres/TEXP-ARITHMETIQUE/exercices/TEXP-ARI-EX-013.tex", "corrige_tex": "chapitres/TEXP-ARITHMETIQUE/corriges/TEXP-ARI-CO-013.tex", "status": "generated"}
% BEGIN-VERIFY
% import math
% def inversibles(n):
%     return [a for a in range(1, n + 1) if any((a * u) % n == 1 for u in range(n))]
% assert inversibles(12) == [1, 5, 7, 11]
% assert len(inversibles(12)) == 4
% assert inversibles(13) == list(range(1, 13))
% assert len(inversibles(13)) == 12
% for n in (12, 13, 20):
%     assert inversibles(n) == [a for a in range(1, n + 1) if math.gcd(a, n) == 1]
% END-VERIFY
\begin{exercice}{TEXP-ARI-EX-013}{3}{18}
Le but est de caractériser les entiers inversibles modulo $n$.
\begin{enumerate}
  \item Montrer que si $a$ est inversible modulo $n$, alors
        $\mathrm{PGCD}(a\,;\,n)=1$.
  \item Réciproquement, montrer que si $\mathrm{PGCD}(a\,;\,n)=1$ alors $a$ est
        inversible modulo $n$.
  \item Dénombrer les entiers de $\{1,\dots,12\}$ inversibles modulo $12$, puis
        modulo $13$. Commenter l'écart.
\end{enumerate}
\end{exercice}
